% ============================================================
%  Carta a los Editores y Árbitros de La Gaceta de la RSME
% ============================================================
\documentclass[a4paper,12pt]{article}
\usepackage[utf8]{inputenc}
\usepackage[T1]{fontenc}
\usepackage[spanish]{babel}
\usepackage{amsmath,amssymb}
\usepackage{geometry}
\geometry{top=3cm, bottom=3cm, left=3.5cm, right=3cm}
\usepackage{parskip}
\setlength{\parskip}{0.75\baselineskip}
\usepackage{hyperref}
\hypersetup{colorlinks=true, urlcolor=blue, linkcolor=black}

\pagestyle{plain}

\begin{document}

\begin{flushright}
  Logroño, noviembre de 2026
\end{flushright}

\medskip
\noindent
Estimado Director,

\medskip
Tenemos el placer de remitirle el artículo

\begin{center}
  \textbf{Dinámica de funciones racionales en la esfera de Riemann:\\
  geometría, cuencas de atracción y visualización}
\end{center}

\noindent
para su consideración en \emph{La Gaceta de la Real Sociedad Matemática
Española}. Aprovechamos esta carta para exponer los aspectos que, a
nuestro juicio, hacen de él una contribución de interés para los
lectores de \emph{La Gaceta}.

\medskip
\noindent\textbf{Contenido matemático y valor divulgativo.}
El artículo aborda la iteración de funciones racionales en la esfera de
Riemann~$\mathbb{P}^1(\mathbb{C})$, uno de los objetos centrales de la
dinámica compleja, con resultados fundacionales debidos a Fatou y Julia
(teoría de conjuntos de Julia y Fatou), Montel (familias normales),
Böttcher (coordenadas normales en entornos de puntos superatractores) y
Sullivan (teorema de no-dominio errante, que completa la clasificación de
componentes de Fatou). La presentación adopta las coordenadas homogéneas
$[z:t]\in\mathbb{P}^1(\mathbb{C})$, que permiten tratar el punto del
infinito en pie de igualdad con los puntos finitos y formular todos los
algoritmos sin casos especiales. Este enfoque ---habitual en geometría
algebraica pero poco frecuente en textos divulgativos de dinámica
compleja--- facilita la comprensión de fenómenos como la permutación de
cuencas de atracción por simetrías globales de la función iterada o la
aparición de cuencas en el infinito. A lo largo del texto se ha buscado
un lenguaje accesible, introduciendo los conceptos mediante expresiones
intuitivas que permitan seguir el razonamiento sin necesidad de un
conocimiento previo especializado en dinámica compleja.

\medskip
\noindent\textbf{Visualización y espacio de color.}
Una aportación específica del artículo es la discusión detallada del
diseño de los <<retratos de cuencas>> desde el punto de vista perceptivo.
Las figuras se generan en el espacio de color CIE~LCHab, derivado del
estándar CIE~1976, que es perceptualmente uniforme: todos los tonos
asignados a las distintas cuencas tienen la misma luminosidad y el mismo
croma, de modo que ninguna cuenca destaca visualmente sobre las demás
por razones ajenas a la dinámica. Esta elección garantiza que el <<retrato
de cuencas>> sea una representación fiel y no sesgada de la partición
dinámica del plano complejo.

\medskip
\noindent\textbf{Accesibilidad e inclusión visual.}
El artículo dedica una sección específica al compromiso con la inclusión
visual en la comunicación matemática. Se discute la paleta Okabe--Ito
(2008), diseñada para ser distinguible por personas con deuteranopía y
protanopía, las formas más comunes de daltonismo, y se incluye una figura
comparativa que ilustra la diferencia entre ambas paletas sobre el mismo
<<retrato de cuencas>>. El paquete \texttt{RiemannBasins.jl} incorpora esta
paleta como opción nativa (\texttt{colorblind\,=\,true}), de modo que
cualquier figura del artículo ---o de trabajos futuros--- puede
regenerarse con garantías de accesibilidad con un único parámetro.
Creemos que este aspecto, aún poco presente en las publicaciones
matemáticas, merece un lugar en un artículo de divulgación.

\medskip
\noindent\textbf{Paquete de software de acceso público.}
El artículo está acompañado por el paquete de código abierto
\texttt{RiemannBasins.jl}, escrito en Julia y disponible en el
repositorio oficial de paquetes de Julia. Los algoritmos implementados utilizan la métrica cordal en~$S^2$ como
criterio de convergencia, y se apoyan en resultados matemáticos
demostrados: la clasificación de las cuencas se basa en el teorema de
las tres componentes de Fatou, y la detección del conjunto de Julia
emplea el criterio de Montel. El paquete
incluye funciones para el <<retrato de cuencas>> en carta afín y sobre la
esfera completa, animaciones sobre familias paramétricas, retratos
tridimensionales sobre~$S^2$, y herramientas para el cálculo de medidas
de las cuencas. Su diseño modular permite extender el análisis a otros
tipos de componentes de Fatou que aparecen en la clasificación de
Sullivan, lo que le otorga una proyección más allá del contenido
específico de este artículo.

\medskip
\noindent\textbf{Adaptación a la edición impresa.}
Dado que \emph{La Gaceta} se publica simultáneamente en versión
impresa ---en blanco y negro--- y en edición digital a color, hemos
añadido algunas indicaciones para los lectores de la versión impresa:
una nota breve al final de la introducción y pies de figura en las
imágenes que dependen del color (Figs.~1, 3, 4, 5, 6, 9 y~10).

\medskip
Esperamos que el artículo pueda resultar de interés para los lectores de
\emph{La Gaceta} y quedamos a disposición del equipo editorial para
cualquier consulta o información adicional que consideren conveniente.

\bigskip
\noindent
Atentamente,

\bigskip\bigskip
\noindent
Víctor Álvarez Aparicio, Luis Javier Hernández Paricio\\
y María Teresa Rivas Rodríguez\\[2pt]
Departamento de Matemáticas y Computación\\
Universidad de La Rioja, Logroño\\[4pt]
\href{mailto:victor.alvarez@unirioja.es}{\texttt{victor.alvarez@unirioja.es}}\\
\href{mailto:luis-javier.hernandez@unirioja.es}{\texttt{luis-javier.hernandez@unirioja.es}}\\
\href{mailto:maria-teresa.rivas@unirioja.es}{\texttt{maria-teresa.rivas@unirioja.es}}

\end{document}
